\scanpage{43}
Pour savoir que $(D_i)_{i\in J}$ y forme bien un
système de ps. droites, il faudra d'abord voir
qu'ils se \unsure{déduisent|relèvent} de la structure de carte
de la façon habituelle, ce qui s'exprime
par la condition (où $s\in S$)
\[\left\{\begin{aligned}
&\text{Si }\rho\in\vec{\unsure{Z|\mathcal Z|\mathfrak Z}}_s,\ \deleted{\unsure{de}}\ \text{et si }{\nu}_s\text{ est le nb de }D_i\text{ passant}\\
&\text{par }s,\ \text{on a }\rho^{\nu_s}(\vec D_i)=(-\vec D_i)\quad\text{pour }\unsure{\text{tte}|\forall}\\
&\text{droite orientée }\inserted{\vec D_i}\text{ passant par }s.
\end{aligned}\right.\tag*{(22)}\]
% slanted
\margin{\uline{NB} Si $\nu_s=2$, alors la condition (22) implique déjà la connaissance de $Z_s\subset{\unsure{\mathfrak S|\mathscr S|\mathcal C}}_{\unsure{\vec D_s|\vec A_s|\vec I_s}}$.}

\deleted{En d'autres termes,}
% slanted
\margin{\uline{\unsure{plus|selon}} la \uline{connexité} de $\mathcal X$ i.e. celle de $K$.}
Il faudra de plus que $\mathcal X$ soit le plan projectif,
ce qui s'exprime par la condition d'EP
\[\underbrace{c_0-c_1}_{1-c(K)}+c_2=1\qquad\text{soit}\quad c_2\deleted{\unsure{(K)}}=c(K)\tag*{(23)}\]
\inserted{\deleted{\unsure{connexité}} degré de connexion de $K$}
% slanted
\margin{\uline{NB} On aura a priori $c_2\leq c(K)$, la formule (23) exprime que la borne sup. est atteinte\ldots}
$c_0=\operatorname{card} S$, $c_1=\operatorname{card} A$ (nb de sommets et
nb d'arêtes) ne dépendent que $K$, tandis que
$c_2$ \inserted{(nb de faces)} dépend bien sûr de $\mathcal X$.

\uline{Construction des} $\vec Z_s$ satisfaisant (22). On
va d'abord définir, pour $\nu_s\geq3$
\[Z_s^{\unsure{\prime|1}}\subset\mathfrak S_{I_s}\tag*{(24)}\]
\unsure{ensem.|comme ens.} de deux ordres circulaires opposés
l'un à l'autre, sur l'ens \inserted{$I_s$} des ps droites passant
% slanted
\margin{i.e. une structure \uline{polygonale} sur cet ens.}
par $s$. Soit $\Delta$ \deleted{$D$} une ps. droite qui ne
passe pas par $s$. \deleted{Alors ch\textsuperscript{une} des ps. droites \inserted{$D_i$}
passant par $s$ est en corr. biunivoque avec l'} \inserted{\unsure{Cons.}, \deleted{\unsure{une $D_s$}} application}
% archive sheet number: 42
